% Stage: — ／ 説明用コード: —
\chapter{エラー訂正とは}\label{ch:qec_intro}

\begin{chaptergoal}
  \begin{itemize}
  \item 古典の誤り訂正（冗長化・多数決）がなぜ量子にそのまま使えないかを理解します。
  \item 測定による状態破壊、No-Cloning 定理、連続的な誤りという、量子特有の 3 つの壁と、パウリ誤りへの離散化を理解します。
  \item データ量子ビットとアンシラを分け、間接測定（シンドローム測定）で誤りを検知する考え方を理解します。
  \item 符号距離 $d$ が訂正能力を決めることを直感レベルで押さえ、後章の合否基準へつなげます。
  \end{itemize}
\end{chaptergoal}

第2章では、FTQC を支える 4 つの制約のひとつとして No-Cloning と間接測定の必要性に軽く触れました。
本章はそこに立ち止まり、\textbf{なぜ}量子の誤り訂正が古典の延長では作れないのかを順に確かめます。
古典の冗長化・多数決から出発し、量子特有の 3 つの壁（測定による破壊、No-Cloning、連続的な誤り）にぶつかり、射影測定による離散化と補助量子ビットを介した間接測定＝シンドローム測定にたどり着く、という一本道を辿ります。
符号そのもの（スタビライザ形式・CSS・表面符号）は次章の主題であり、本章は「なぜその形にならざるを得ないか」という動機づけに専念します。

\section{古典と量子の誤り訂正}\label{sec:classical_vs_quantum}

\subsection{冗長化と多数決}

古典の通信路や記憶装置でも、ビットは確率的に反転します。
最も素朴な対策は\textbf{冗長化}です。
1 ビットの情報 $b\in\{0,1\}$ を 3 回コピーして $bbb$（$000$ または $111$）として送り、受信側は多数決で復号します。

たとえば $111$ を送って 1 ビットだけが反転し $101$ になったとしても、多数決は $1$ を返し、元の情報は正しく復元されます。
各ビットが独立に確率 $p$ で反転するとき、多数決が誤るのは 2 ビット以上が反転した場合に限られるので、誤り率はおよそ $p$ から $3p^2$ へ下がります（$p$ が十分小さいとき）。
コピー数（符号長）を増やすほど、この誤り率はさらに下がります。
この「冗長化 → 直接読み出し → 多数決」という手順が、量子ではそのまま使えません。

\subsection{量子でそのまま使えない理由}

古典の手順は、実は 3 つのことを暗黙に仮定しています。

\begin{itemize}
  \item 情報を複製できる（$b$ を 3 つのビットにコピーできる）。
  \item 各コピーを直接読み出しても、情報そのものは壊れない（読んだ後も同じ $b$ が残っている）。
  \item 生じる誤りは $0 \leftrightarrow 1$ の反転（離散的な反転）のみである。
\end{itemize}

量子ビットの重ね合わせ（第1章）を守ろうとすると、これら 3 つの仮定はいずれも成り立ちません。
次節で、この 3 つの壁とパウリ誤りへの離散化を順に確認します。

\section{量子特有の壁}\label{sec:quantum_walls}

\subsection{測定による状態破壊}

第1章で見たとおり、計算基底での測定は状態を対応する基底へ\textbf{射影}し、重ね合わせを壊します。
一般の状態 $\ket{\psi}=\alpha\ket{0}+\beta\ket{1}$ を測定してしまうと、結果が $0$ か $1$ のどちらかに定まり、係数 $\alpha,\beta$ という情報自体が失われます。

古典の多数決は「各コピーを読んで比較する」手順でした。
量子データ量子ビットに対してこれを文字通り行うと、比較のたびに重ね合わせを壊してしまい、守りたかった量子情報そのものが消えてしまいます。
つまり\textbf{データ量子ビットを直接測定する}という古典の発想は、量子では誤り訂正の対象を破壊する自己矛盾を起こします。

\subsection{No-Cloning}

もう一つの仮定である「複製」も、量子では一般に禁止されています。
これを\textbf{No-Cloning 定理}と呼びます。

任意の未知状態 $\ket{\psi}$ に対して
\begin{equation}
  U\bigl(\ket{\psi}\otimes\ket{0}\bigr) = \ket{\psi}\otimes\ket{\psi}
\end{equation}
を満たす、状態によらない\textbf{単一の}ユニタリ $U$ が存在すると仮定してみます。
計算基底では
\begin{equation}
  U\ket{00} = \ket{00},
  \qquad
  U\ket{10} = \ket{11}
\end{equation}
となるはずです。$U$ はユニタリなので線形性を持ち、重ね合わせ $\ket{+}=\frac{1}{\sqrt2}(\ket0+\ket1)$ に対しては
\begin{equation}
  U\bigl(\ket{+}\otimes\ket{0}\bigr)
  = \frac{1}{\sqrt2}\bigl(U\ket{00}+U\ket{10}\bigr)
  = \frac{1}{\sqrt2}\bigl(\ket{00}+\ket{11}\bigr)
  = \ket{\Phi^+}
\end{equation}
となります。これは第1章のベル状態そのものです。
一方、$U$ が本当にクローニングを行うなら
\begin{equation}
  U\bigl(\ket{+}\otimes\ket{0}\bigr) \overset{!}{=} \ket{+}\otimes\ket{+}
  = \frac12\bigl(\ket{00}+\ket{01}+\ket{10}+\ket{11}\bigr)
\end{equation}
にならなければなりません。
ベル状態 $\ket{\Phi^+}$ と積状態 $\ket{+}\otimes\ket{+}$ は異なる状態なので、$0,1$ をクローンできる $U$ は $\ket{+}$ をクローンできず、矛盾します。
したがって、任意の未知量子状態に対して働く万能なクローニング装置は存在しません。

\begin{notebox}
  ここで禁止されるのは「\textbf{未知}かつ\textbf{任意}の状態を非破壊に複製すること」です。
  既知の直交状態（たとえば「$\ket0$ か $\ket1$ のどちらかだと分かっている」）ならコピーできますし、量子もつれを使ったテレポーテーションのように状態を\textbf{移す}ことも別問題です。
  誤り訂正で問題になるのは、計算途中の未知の重ね合わせ係数 $\alpha,\beta$ を保ったまま増やせない、という点です。
\end{notebox}

\subsection{連続的な誤り}

古典ビットの誤りは $0 \leftrightarrow 1$ の反転だけです。
しかし量子ビットはブロッホ球上の任意の重ね合わせ状態をとるため、微小な回転
\begin{equation}
  R_x(\theta) = \ee^{-\ii \theta X/2}, \qquad R_z(\phi) = \ee^{-\ii \phi Z/2}
\end{equation}
のように、連続的で無限通りの誤りが起こり得ます。
連続的な回転角 $\theta$ や位相のずれ $\phi$ を無限の精度で特定して元に戻すことは、一見すると不可能なように思えます。

\subsection{パウリ誤りへの離散化}

この「連続的な誤り」という壁は、量子測定の射影という性質を逆に利用して解決します。
これを\textbf{誤りの離散化}と呼びます。

第1章で定義したパウリ行列 $\{I, X, Y, Z\}$ を思い出してください。
任意の 1 量子ビット演算子（誤り演算子）$E$ は、パウリ行列の複素線形結合として
\begin{equation}
  E = c_0 I + c_1 X + c_2 Y + c_3 Z \qquad (c_i \in \mathbb{C})
\end{equation}
と展開できます。
ここで、$X$ はビット反転（古典の反転）、$Z$ は位相反転（$\ket{\pm}$ 基底での反転）、そして $Y = \ii XZ$ はビット反転と位相反転の両方が起きた誤りに対応します。

符号化された状態 $\ket{\psi}$ が誤り $E$ を受けたとき、$E\ket{\psi}$ に対して後述するシンドローム測定（パウリ演算子の固有値を測る射影測定）を行うと、状態は
\begin{equation}
  I\ket{\psi} \ (\text{誤りなし}), \quad X\ket{\psi} \ (\text{ビット反転}), \quad Z\ket{\psi} \ (\text{位相反転}), \quad Y\ket{\psi} \ (\text{両方})
\end{equation}
のいずれか 1 つに確率的に確定します（ここでは、符号がこれらの誤りを互いに異なるシンドロームで区別できるものとします）。
つまり、測定する前は連続的だった誤りが、測定によって離散的なパウリ誤りのどれか 1 つに確定します。

たとえば、微小な $X$ 回転誤り $R_x(\theta) = \cos(\theta/2) I - \ii \sin(\theta/2) X$ が起きたとします。
この誤りが符号化された量子ビットの 1 つに起きたとしてシンドローム測定を行うと、状態は確率 $\cos^2(\theta/2)$ で「誤りなし（$I$）」に、確率 $\sin^2(\theta/2)$ で「$X$ 誤り」に射影されます。
ひとたびどちらかに射影されれば、シンドロームの結果を見て「何もしない」か「$X$ をかけて直す」かのどちらかを行うだけで、元の状態に完全に復元できます。

したがって、量子誤り訂正では\textbf{ビット反転 $X$ と位相反転 $Z$ の 2 種類（およびその積 $Y$）さえ正しく検出・訂正できれば、すべての連続的な誤りに対処できる}ことになります。

\subsection{間接測定の必然性}

ここまでに挙げた 3 つの壁と離散化の性質を整理すると、次の要件が浮かび上がります。

\begin{itemize}
  \item データ量子ビットを直接測定してはいけない（重ね合わせが壊れる）。
  \item データ量子ビットを複製してもいけない（未知状態は複製できない）。
  \item 誤りを連続的なまま無限精度で測定する必要はなく、パウリ誤りへ離散化して検出すれば有限種類の訂正で済む。
\end{itemize}

データそのものを壊さずに、かつ複製も使わずに、「どのパウリ誤りがどこで起きたか」だけを知るにはどうすればよいでしょうか。
残された道は、データ量子ビットそのものではなく、データに\textbf{別の量子ビットを量子もつれで絡ませて}、その別の量子ビット越しに誤りの兆候だけを読み出す、という間接的な方法です。
次節で、この役割を担う補助量子ビット（アンシラ）とシンドローム測定を具体的に見ます。

\section{補助量子ビットとシンドローム測定}\label{sec:syndrome}

\subsection{データ vs アンシラ}

量子誤り訂正では、量子ビットを 2 種類に分けます。

\begin{itemize}
  \item \textbf{データ量子ビット}：守りたい論理情報（重ね合わせ係数）を保持します。直接測定しません。
  \item \textbf{補助量子ビット（アンシラ）}：誤りの兆候を運ぶためだけに使い捨てで用意し、測定してよいのはこちらです。
\end{itemize}

アンシラをデータと量子もつれさせ、アンシラだけを測定することで、データの重ね合わせを壊さずに「誤りが起きたかどうか」という情報だけを取り出します。
この一連の操作を\textbf{シンドローム測定}と呼びます。

\subsection{パリティ検査の直感}

具体例として、3 個のデータ量子ビットで論理情報を
\begin{equation}
  \ket{\psi_L} = \alpha\ket{000} + \beta\ket{111}
\end{equation}
のように冗長化する場合を考えます（古典の $bbb$ の量子版に相当します）。
ここで注意したいのは、この状態は未知の量子状態 $\alpha\ket0+\beta\ket1$ を 3 つに複製した積状態 $(\alpha\ket0+\beta\ket1)^{\otimes 3}$ ではなく、3 量子ビットの状態空間に埋め込んだ量子もつれ状態である点です。
したがって、No-Cloning 定理とは矛盾しません。

このとき、隣接ビットの偶奇（パリティ）を表す 2 つの演算子
\begin{equation}
  S_1 = Z_0 Z_1, \qquad S_2 = Z_1 Z_2
\end{equation}
を測定すれば、係数 $\alpha,\beta$ を読まずに「反転が起きたかどうか」だけが分かります。
$\ket{000}$ と $\ket{111}$ はどちらも $Z_0Z_1=+1$ かつ $Z_1Z_2=+1$ の固有状態なので、これらのパリティ測定は論理情報 $\alpha,\beta$ を壊しません。

パリティの測定は、図~\ref{fig:parity_check} のように $\ket0$ に初期化した 2 個の補助量子ビット（アンシラ）$a_1, a_2$ を用意し、CNOT でデータの偶奇をアンシラへ写してからアンシラだけを測定することで行います。

\begin{figure}[htbp]
  \centering
  \begin{quantikz}
    \lstick{$q_0$} & \ctrl{3} & \qw & \qw & \qw & \qw \\
    \lstick{$q_1$} & \qw & \ctrl{2} & \ctrl{3} & \qw & \qw \\
    \lstick{$q_2$} & \qw & \qw & \qw & \ctrl{2} & \qw \\
    \lstick{$a_1:\ket{0}$} & \targ{} & \targ{} & \qw & \qw & \meter{} \\
    \lstick{$a_2:\ket{0}$} & \qw & \qw & \targ{} & \targ{} & \meter{}
  \end{quantikz}
  \caption{3 量子ビット反復符号のシンドローム測定回路。データ $q_0, q_1$ の偶奇（パリティ $Z_0Z_1$）をアンシラ $a_1$ に、データ $q_1, q_2$ の偶奇（パリティ $Z_1Z_2$）をアンシラ $a_2$ に CNOT で写して測定します。データの重ね合わせを壊さずに誤りの有無と位置を特定できます。}
  \label{fig:parity_check}
\end{figure}

アンシラ $a_1, a_2$ の測定結果を $(s_1, s_2) \in \{0, 1\}^2$ と表します（固有値 $+1$ を $0$、$-1$ を $1$ に対応させます）。
この測定値の組を\textbf{シンドローム}と呼びます。

あるデータ量子ビットに $X$ 誤りが起きると、パウリ演算子の反可換性 $ZX = -XZ$ により、その量子ビットを含むパリティ演算子の固有値だけが $-1$ に反転します。
たとえば $q_0$ に $X$ 誤りが起きると、$Z_0$ と反可換である $S_1 = Z_0 Z_1$ の符号が反転して $s_1=1$ となりますが、$q_0$ を含まない $S_2 = Z_1 Z_2$ は可換なため $s_2=0$ のままです。
表~\ref{tab:ch3_syndrome} にシンドロームと推定される誤り、および訂正操作の対応を示します。

\begin{table}[htbp]
  \centering
  \caption{3 量子ビット反復符号のシンドローム表}
  \label{tab:ch3_syndrome}
  \begin{tabular}{cccl}
    \hline
    $s_1 \, (Z_0Z_1)$ & $s_2 \, (Z_1Z_2)$ & 推定される誤り & 訂正操作 \\
    \hline
    $0$ & $0$ & なし & 何もしない \\
    $1$ & $0$ & $q_0$ に $X$ & $q_0$ に $X$ をかける \\
    $1$ & $1$ & $q_1$ に $X$ & $q_1$ に $X$ をかける \\
    $0$ & $1$ & $q_2$ に $X$ & $q_2$ に $X$ をかける \\
    \hline
  \end{tabular}
\end{table}

シンドロームから誤りの位置が特定できれば、その量子ビットにもう一度 $X$ を作用させることで、$X^2 = I$ により元の状態に戻せます。
重要なのは、シンドローム測定が「$\alpha\ket{000}+\beta\ket{111}$ のどちらであるか」を一切読まず、「反転が起きたかどうかとその位置」だけを教えてくれる点です。
だからこそ、データの重ね合わせを壊さずに誤りを検知・訂正できます。
これが、古典の「直接読んで多数決」を、量子で「量子もつれで間接的に読む」へと置き換えた形です。
第5章では、この手順を周期的に繰り返す\textbf{符号サイクル}として時間軸に沿って再整理します。

\section{符号距離の壁}\label{sec:distance_wall}

前節で導入した 3 量子ビット反復符号は、間接測定の仕組みを理解する絶好の例題ですが、実用上は 2 つの明確な限界を持っています。

第 1 の限界は、\textbf{2 か所以上のビット反転誤りを訂正できない}ことです。
たとえば $q_0$ と $q_1$ に同時に $X$ 誤りが起きると、状態は $\alpha\ket{110} + \beta\ket{001}$ に変化します。
このときのシンドロームを計算すると $(s_1, s_2) = (0, 1)$ となります。
復号器は表~\ref{tab:ch3_syndrome} に従って「$q_2$ に $X$ 誤りが起きた」と判断し、$q_2$ に $X$ をかけて訂正を試みます。
その結果、状態は $\alpha\ket{111} + \beta\ket{000}$ となり、論理ビットが反転（論理誤り）して訂正に失敗してしまいます。

第 2 の限界は、\textbf{位相反転誤り（$Z$）をまったく検出できない}ことです。
たとえば $q_0$ に $Z$ 誤りが起きた場合、論理状態は $\alpha\ket{000} - \beta\ket{111}$ に変わります。
しかし $Z_0$ はパリティ演算子 $S_1 = Z_0 Z_1$ および $S_2 = Z_1 Z_2$ の両方と可換であるため、シンドロームは $(0, 0)$ のままであり、誤りが起きたことにすら気づけません。

したがって、本格的な量子誤り訂正を実現するには、ビット反転 $X$ だけでなく位相反転 $Z$ も同時に防ぐことができ、かつより多くの誤りに耐えられる符号を設計しなければなりません。

符号がいくつまでの誤りを見分けられるかは、一般に\textbf{符号距離} $d$ という 1 つの整数で決まります。
符号距離 $d$ は「非自明な論理演算子の最小重み」（＝検出されずに論理状態を別の論理状態へ変えてしまう最小の誤りの個数）として定義されます。
この符号距離 $d$ により、重み $d-1$ 以下のパウリ誤りを検出できます。また、訂正半径を
\begin{equation}
  t = \left\lfloor \frac{d-1}{2} \right\rfloor
\end{equation}
とおくと、重み $t$ 以下の任意のパウリ誤りを一意に訂正できます。
3 ビット反復符号はビット反転に対して $d=3$ に相当し、重み $d-1 = 2$ 以下のビット反転を検出でき、$\lfloor (3-1)/2 \rfloor = 1$ 個までのビット反転を一意に訂正できる、という先ほどの観察と整合します。

$d$ を大きくすれば訂正能力は上がりますが、その分だけ物理量子ビット数・アンシラ数・シンドローム測定の回数が増えます（第2章の符号距離の制約を参照）。
ビット反転 $X$ と位相反転 $Z$ の双方を同時に保護し、高い符号距離 $d$ を実現する符号がどのように構成され、論理演算子やスタビライザとしてどう記述されるかは次章の主題です。
また、本書のシミュレータでは合否基準の中心に「回路から求めた最短の論理エラー（$\mathrm{shortest\_graphlike\_error}$）が符号距離 $d$ と一致するか」を据えます。これは、実装した符号サイクルが\textbf{意図した距離をきちんと実現できているか}を機械的に確認するためであり、詳細は第8章で定義します。

\section{やること・スコープ外}\label{sec:qec_intro_scope}

\subsection{やること}

\begin{itemize}
  \item 古典の冗長化・多数決が量子でそのまま使えない理由（測定による破壊、No-Cloning、連続的な誤り）とパウリ誤りへの離散化を示す。
  \item データ量子ビットとアンシラを分け、シンドローム測定という間接測定の考え方を導入する。
  \item 符号距離 $d$ と訂正能力の関係を、3 ビット反復符号の限界と合わせて直感的に示す。
\end{itemize}

\subsection{スコープ外}

\begin{itemize}
  \item スタビライザ形式による一般的な符号の定義、CSS 条件、表面符号の幾何（第4章）。
  \item シンドロームからの実際の復号アルゴリズム（MWPM、BP+OSD 等、第8章以降）。
  \item 符号サイクルの時間発展、パウリフレーム追跡との統合（第5章以降）。
\end{itemize}

\section{まとめ}\label{ch3:summary}

\begin{itemize}
  \item 古典の誤り訂正は「複製して直接読み、多数決を取る」手順に依存します。
  \item 量子には測定による状態破壊、No-Cloning 定理、連続的な誤りという 3 つの壁がありますが、射影測定によるパウリ誤りへの離散化と、補助量子ビットによる間接測定（シンドローム測定）によって克服できます。
  \item 3 量子ビット反復符号では、データとアンシラの間の CNOT でパリティ（$Z_0Z_1, Z_1Z_2$）を測ることで、論理情報を壊さずにビット反転 $X$ を検知・訂正できます。
  \item しかし反復符号は 2 個以上のビット反転や位相反転誤り（$Z$）には対処できません。$X$ と $Z$ の両方を訂正できる本格的な符号が必要です。
  \item 符号距離 $d$ が誤りの検出可能重み（$d-1$ 以下）および訂正半径 $t = \lfloor (d-1)/2 \rfloor$ を決め、$\mathrm{shortest\_graphlike\_error}=d$ という形で後章の合否基準につながります。
\end{itemize}

次章では、本章で得た「アンシラでパリティを測る」という直感を、スタビライザ形式・CSS 符号・表面符号という一般的な枠組みに広げます。

\begin{exercise}
  \begin{enumerate}
    \item 古典の 3 ビット反復符号における多数決が、量子のデータ量子ビットに対してそのまま使えない理由を、測定による破壊と No-Cloning のそれぞれの観点から一文ずつ説明しなさい。
    \item 本章の No-Cloning の証明で、$U\ket{00}=\ket{00}$、$U\ket{10}=\ket{11}$ を仮定した場合に $U\ket{01}$、$U\ket{11}$ がどうなるべきかを考え、$\ket{-}=\frac{1}{\sqrt2}(\ket0-\ket1)$ についても同様の矛盾が起きることを確かめなさい。
    \item $\ket{\psi_L}=\alpha\ket{000}+\beta\ket{111}$ の $q_0$ に $X$ 誤りが起きたとき、$Z_0Z_1$ と $Z_1Z_2$ の測定値（シンドローム）がそれぞれどうなるかを求め、誤りの位置を特定できることを確かめなさい。
    \item 3 量子ビット反復符号の状態 $\ket{\psi_L}=\alpha\ket{000}+\beta\ket{111}$ において、$q_0$ に位相反転誤り $Z$ が生じた場合を考えます。パリティ測定 $Z_0Z_1$ および $Z_1Z_2$ のシンドローム $(s_1, s_2)$ がどうなるかを計算し、なぜこの符号では位相反転誤りを検出できないのかを説明しなさい。
  \end{enumerate}
\end{exercise}
